\documentclass[10pt,a4paper]{ctexart}
\usepackage[margin=1.4cm]{geometry}
\usepackage{booktabs}
\usepackage{enumitem}
\usepackage{amsmath,amssymb}
\usepackage{graphicx}
\setlist{nosep,leftmargin=1.2em,topsep=1pt}
\setlength{\parskip}{1pt}
\setlength{\parindent}{0pt}
\renewcommand{\arraystretch}{0.95}
\pagestyle{empty}
\usepackage{titlesec}
\titlespacing*{\section}{0pt}{4pt}{2pt}
\titleformat{\section}{\normalsize\bfseries}{\thesection.}{0.5em}{}

\begin{document}
\begin{center}\large\bfseries 赛道一：ModelNet40 三维点云分类 —— 设计思路（PointMLP）\end{center}
\vspace{-0.6em}

\section{任务与目标}
输入：每个样本 10000 个点，每点 6 维 $(x,y,z,n_x,n_y,n_z)$。\\
输出：40 类形状标签。\\
评分目标：实例准确率 (iAcc) $\geq 92\%$，类平均准确率 (cAcc) $\geq 90\%$。

\section{核心思路：为什么选 PointMLP}

原方案用 5 个模型（PointNet2~SSG/Msg + DGCNN）集成来凑够分数，但单模型几乎全部不达标。我们希望能用一个模型独立过线，于是调研了近年 SOTA：

\begin{center}\small
\begin{tabular}{lcll}
\toprule
方法 & 会议 & OA & 特点 \\
\midrule
PointNet++ & NeurIPS'17 & 91.9\% & 分层点集抽象，FPS+ball query \\
DGCNN & TOG'19 & 92.9\% & 动态图 CNN，EdgeConv  \\
Point Transformer & ICCV'21 & 93.7\% & 向量自注意力 \\
\textbf{PointMLP} & \textbf{ICLR'22} & \textbf{94.1\%} & \textbf{纯 MLP，无卷积/注意力} \\
PointNeXt & NeurIPS'22 & 94.1\% & PointNet++改进训练 \\
Point-MAE & ECCV'22 & 93.8\% & 掩码预训练（需额外数据） \\
\bottomrule
\end{tabular}
\end{center}

\textbf{选择 PointMLP 的理由}：
\begin{enumerate}
\item 纯 MLP 架构，不依赖卷积或注意力，计算高效；
\item 官方 OA 94.1\% / mAcc 91.5\%，单模型即可过线；
\item 有完整开源代码，可直接复现；
\item xyz-only 输入（无需法线），减少了跨分布的泛化风险。
\end{enumerate}

\section{PointMLP 架构详解}

PointMLP 的核心洞察是：\textbf{点云分类的关键不是复杂的算子，而是合适的局部几何建模}。架构由三部分组成：

\subsection{LocalGrouper（局部聚合）}
\begin{itemize}
\item FPS（最远点采样）在空间中均匀降采样中心点；
\item KNN 在每个中心点周围取 $k=24$ 个最近邻；
\item \textbf{Geometric Affine 归一化}：以锚点（中心点）为参考，将邻域内各点标准化：$\mathbf{x}' = (\mathbf{x} - \mathbf{x}_{\text{anchor}}) / \sigma$，再通过可学习参数 $\alpha,\beta$ 做仿射变换：$\mathbf{y} = \alpha \cdot \mathbf{x}' + \beta$；
\item 最后拼接锚点特征作为残差上下文（类似残差连接的隐式形式）。
\end{itemize}

\subsection{Pre \& Post Extraction（前后残差 MLP）}
\begin{itemize}
\item \textbf{PreExtraction}：在邻域的 $k$ 个点上做共享 MLP + max-pool，升维到目标通道数，再加若干个残差 block；
\item \textbf{PosExtraction}：在池化后的全局特征上做通道间残差 MLP。
\item 双层残差设计（邻域内 + 通道间）保证深层梯度流动。
\end{itemize}

\subsection{整体结构}
4 级层级结构，通道逐级翻倍（$64\to128\to256\to512\to1024$）、点数逐级减半（$1024\to512\to256\to128\to64$）：

\begin{center}\small
\begin{tabular}{cccccc}
\toprule
Stage & 输入点数 & KNN $k$ & 输入通道 & 输出通道 & 残差 block 数 \\
\midrule
Embedding & 1024 & ---  & 3    & 64  & 1 (Conv1d) \\
Stage 1  & 512  & 24 & 64   & 128 & 2+2 \\
Stage 2  & 256  & 24 & 128  & 256 & 2+2 \\
Stage 3  & 128  & 24 & 256  & 512 & 2+2 \\
Stage 4  & 64   & 24 & 512  & 1024& 2+2 \\
Classifier & --- & --- & 1024 & 40 & FC$\to$512$\to$256$\to$40 \\
\bottomrule
\end{tabular}
\end{center}

全局 max-pool 聚合后将 $1024$ 维特征送入三层 FC 分类头，中间带 BN+ReLU+Dropout(0.5)。总参数量 \textbf{13.27M}。

\section{训练过程}

\subsection{数据预处理（\texttt{preprocess.py}）}
将教师提供的 $\sim$9843 个 \texttt{.txt} 文件（每文件 10000 行 × 6 列）读入并统一为 \texttt{data.npy}（$\sim$1.2\,GB, fp16），同时生成 \texttt{labels.npy} 和 \texttt{shape\_names.npy}。后续通过 mmap 读取，避免大量小文件 I/O。

\subsection{分层验证划分}
\texttt{stratified\_split(labels, val\_ratio=0.1, seed=2026)} 对每类独立做 9:1 分层采样，得到固定的 8877 训练 / 966 验证样本。所有超参调优仅看 val，test 最后一次运行。

\subsection{训练配方（PointMLP 官方配方）}
\begin{itemize}
\item \textbf{优化器}：SGD momentum=0.9, weight\_decay=$2\times10^{-4}$（不用 Adam/AdamW；PointMLP 论文指出 SGD 是该架构的最佳搭档）；
\item \textbf{学习率}：Cosine annealing，从 0.1 退火到 0.005；
\item \textbf{损失函数}：CrossEntropyLoss(label\_smoothing=0.2)；
\item \textbf{混合精度}：\texttt{torch.cuda.amp} + GradScaler，训练加速约 30\%；
\item \textbf{Batch size}：32；
\item \textbf{训练轮数}：200 epochs，每 2 个 epoch 在 val 上评测一次；
\item \textbf{增强}（仅 xyz）：
  \begin{itemize}
  \item 各向异性缩放 $[0.66, 1.5]^3$
  \item 平移 $\pm 0.2$
  \item 不做旋转（xyz-only 模型旋转增强反而有害）
  \end{itemize}
\end{itemize}

训练在单张 RTX A6000 (49GB) 上完成，约 97 秒/epoch，总计 \textbf{$\sim$5.4 小时}。

\subsection{Checkpoint 策略}
每轮保存三个 ckpt：
\begin{itemize}
\item \texttt{pointmlp.pt} —— 至今 val iAcc 最高者；
\item \texttt{pointmlp\_bestcacc.pt} —— 至今 val cAcc 最高者；
\item \texttt{pointmlp.pt.last} —— 最后一轮。
\end{itemize}

\section{TTA 推理（\texttt{predict.py}）}
\begin{itemize}
\item 对每个样本，先在整个 10000 点上做 center + unit-ball 归一化；
\item 随机子采样 1024 个点，过模型取 softmax，重复 10 次；
\item 10 轮 softmax 求平均后 argmax 得最终预测。
\end{itemize}
支持三种输入模式：npy 文件、txt 目录、id 列表+根目录。另设 \texttt{--fast} 应急模式（TTA=1，CPU 可跑 $\sim$30s）。

\section{实验结果}

\subsection{验证集（966 样本，seed=2026 val split）}

\begin{center}\small
\begin{tabular}{lccc}
\toprule
Checkpoint & Epoch & val iAcc & val cAcc \\
\midrule
best iAcc & 134 & \textbf{93.48\%} & 90.00\% \\
best cAcc & 182 & 93.37\% & \textbf{90.61\%} \\
last      & 200 & 92.96\% & 88.78\% \\
\bottomrule
\end{tabular}
\end{center}

\subsection{测试集（2468 样本，教师下发）}

采用 best\_cAcc checkpoint 进行最终评测：

\begin{center}\small
\begin{tabular}{lccc}
\toprule
方案 & test iAcc & test cAcc & 是否达标 \\
\midrule
旧方案（5 模型集成 + TTA）& 92.38\% & 89.39\% & $\times$ \\
\textbf{PointMLP (本方案)} & \textbf{93.19\%} & \textbf{90.34\%} & \textbf{$\checkmark$} \\
\bottomrule
\end{tabular}
\end{center}

\textbf{双指标均超过 92\% / 90\% 门槛。} 单 PointMLP 模型在 iAcc 上比旧 5 模型集成提升 0.81\%，cAcc 提升 0.95\%。

\subsection{主要困难类别}

flower\_pot 在所有方案中持续表现最差（test 上仅 4/20 recall），系花瓶/花盆/盆栽三类几何形状高度相似导致的固有分类歧义。ICML 2023 论文 \textit{``ModelNet40: What is left?''} 也指出该问题属于标签噪声和类间固有歧义，非模型能力问题。

\section{对比旧方案}

\begin{center}\small
\begin{tabular}{lccc}
\toprule
维度 & 旧方案（5 模型集成） & 新方案（PointMLP） \\
\midrule
模型数量 & 5 个（3 类架构） & 1 个 \\
总参数量 & $\sim$8.2M × 5 & 13.27M \\
训练时间 & $\sim$7--8 h $\times$ 5 & $\sim$5.4 h \\
推理时间 (TTA=10) & $\sim$10--20 min & $\sim$5--10 min \\
单模 iAcc/cAcc & 全部 $<92\%$/$<90\%$ & \textbf{93.19\%}/\textbf{90.34\%} \\
部署复杂度 & 需同时加载 5 个模型 & 单文件 ckpt \\
\bottomrule
\end{tabular}
\end{center}

\section{结论}

PointMLP 以更简单的架构（纯 MLP，无卷积/注意力）、更少的模型数量（1 vs 5）、更短的训练时间，在 ModelNet40 上实现了超越原集成方案的结果，且无需法线信息（仅 xyz 输入），降低了跨预处理流水线的分布偏移风险。

\end{document}
